\section{Residues at primitive-recursively named singularities}
\label{sec:namedresidue}

The residue theorem in the repository takes its singularity centres in
$\mathbb Q(i)$. Its extension to rectifiable curves retains that restriction.
A pole of a rational function with rational coefficients, for example
$+\sqrt{2}$ for $1/(z^2-2)$, need not have rational coordinates.
There are two different uses of the centre in the proof. Laurent
coefficients use it as an origin, while the planar filling algorithm
compares coordinates to draw a mesh. We translate the analytic origin
by a name and draw the mesh around one rational sample. A strip joining
the two contours proves that the computed residue is the coefficient at
the intended origin.

\paragraph{The rational case used below.}
Write $Q_t(c)=\{z:\lVert z-c\rVert_\infty\le t\}$ and put
$|z|_1=|\Re z|+|\Im z|$. Norms in error estimates below are $|\cdot|_1$.
Two results from the repository's
\href{https://github.com/VladimirReshetnikov/ProveIt/blob/5f3be02388442f16487d2dc5c2529961dd783d2b/Logic/PrimitiveRecursiveArithmetic/docs/proofs/analysis/cauchy_residue_ipra/report.tex}{Cauchy and residue report}
are used as mathematical lemmas. First, the integral on the boundary of
a finite rational triangle chain is zero under a uniform first-order
remainder estimate on its carrier. Second, on the rational annulus
$Q_{8r}(0)\setminus\operatorname{int} Q_{r/8}(0)$, the same effective
holomorphy data give Laurent coefficients $a_m$, a unique expansion on
$\partial Q_r(0)$, and $\int_{\partial Q_r(0)}g=2\pi i a_{-1}$.
The annular result derives the coefficients by expanding the two separated
Cauchy kernels on the outer and inner boundary. Both finite geometric
remainders have ratio at most $1/2$. For inner and outer magnitude bounds
at most $M$, the $2N$ terms with indices $-N,\ldots,N-1$ have uniform
error at most $40M2^{-N}$. The moment integral of each integer power
vanishes except for $z^{-1}$, whose integral is $2\pi i$. The linked
report gives the human proofs in its sections on the annular Cauchy
formula and the final residue error budget.

\subsection{Translating a holomorphic chart}

A name $A(n)\in\mathbb Q(i)$ for $a$ is a fixed PR function satisfying
$|A(m)-A(n)|_1\le 2^{-m}+2^{-n}$. A holomorphic chart on a rational
carrier supplies PR evaluators $F(z,n),D(z,n)$, each coherent with
error $2^{-n}$, magnitude bounds $M,B$, and nondecreasing PR moduli
$\nu,\mu$. In name notation the conditions at rational points $z,w$ are
\begin{align*}
 |w-z|_1\le 2^{-\nu(k)}
   &\ \Longrightarrow\ |f(w)-f(z)|_1\le 2^{-k},\\
 |w-z|_1\le 2^{-\mu(k)}
   &\ \Longrightarrow\ |f(w)-f(z)-f'(z)(w-z)|_1
                      \le 2^{-k}|w-z|_1 .
\end{align*}
The first bound at sample precisions $n,m$ has extra slack $2^{-n}+2^{-m}$.
The second at a common precision $n$ has extra slack
$2^{-n}(2+|w-z|_1)$. These clauses, with the two coherence and magnitude
bounds, are first-order statements on rational codes.

\paragraph{Lemma (a modulus for the derivative).}
On that carrier $\mu_D(k)=\mu(k+2)$ works for $f'$. Indeed, let
$h=w-z\ne0$, $d=|h|_1$. Add the two first-order remainders, one in
each direction, and obtain
$|(f'(z)-f'(w))h|_1\le 2\cdot 2^{-(k+2)}d$ when
$d\le 2^{-\mu(k+2)}$. For $h=x+iy$ one has
$|h^{-1}|_1=d/(x^2+y^2)\le 2/d$. Submultiplicativity gives the
claimed $2^{-k}$ bound. The zero-displacement case uses coherence at a
single rational point. Division in this argument is a decidable
operation on nonzero rational $h$. At precision $m$ the two remainders
contribute an additional $4(2+d)2^{-m}/d$ after division. This factor
can be budgeted for each rational pair before taking the arbitrary-error
consequence. For example, to make the extra term at most $2^{-v}$, take
$m=v+\sigma(4(2+d)/d)$, where throughout this section
$\sigma(u)$ is the least $s\ge0$ with $2^s\ge\max(1,u)$.
For a nonnegative fraction $p/q$, $q\ge1$, the search is bounded by $p+1$.

\paragraph{Lemma (translation).}
Choose $r>0$ rational, $b$ with $2^{-b}\le r/64$, and
$c=A(b)$. Assume the chart clauses on the rational closed shell
\[
 S(c,r)=\{z:r/16\le\lVert z-c\rVert_\infty\le 9r\}.
\]
For rational $x$ with $r/8\le\lVert x\rVert_\infty\le8r$, the
translated function $g(x)=f(a+x)$ has PR holomorphy data on this
fixed rational annulus. In particular the preceding rational Laurent
theorem applies to $g$.

To construct the data, observe that for every $t\ge b$,
$\lVert A(t)-c\rVert_\infty\le r/32$. Thus $x+A(t)$ belongs to $S(c,r)$
because $r/8-r/32\ge r/16$ and $8r+r/32\le9r$. To produce the value
and derivative samples at precision $n$, use
\[
 t(n)=\max\{b,1+\nu(n+2),1+\mu(n+4)\}
\]
and evaluate $F(x+A(t(n)),n+2)$ and $D(x+A(t(n)),n+2)$.
Every later centre sample is within $2^{1-t(n)}$ of $A(t(n))$.
The input modulus and the derivative lemma make the centre error at
most $2^{-(n+2)}$ in either value. The evaluator contributes at most
another $2^{-(n+2)}$. In fully rational form, compare two output
samples by inserting a common later centre and a common fine evaluator
precision, then discard the arbitrarily small positive slack by
cross-multiplication. Thus the displayed programs each have error
at most $2^{-n}$ and the required coherence. A translated modulus or
first-order bound for $x,y$ follows by evaluating both at the \emph{same}
later centre $A(T)$. Their displacement is exactly $y-x$. Taking
arbitrary rational precision recovers $\nu,\mu$ unchanged. Replacing
the three resulting names by their output samples introduces precisely
the allowed evaluator slack. Magnitude bounds transport in the same
way. If two input names represent the same centre, the common-sample
comparison with $\nu$ and $\mu_D$ gives equal translated names wherever
the indicated shells are certified.

\subsection{The moving-centre residue theorem}

Let $\Gamma$ be a closed rational polygonal chain and let $A_j(n)$
name finitely many centres $a_j$ (jointly PR in $j,n$ if their number is
a parameter). For each $j$ choose $r_j,b_j,c_j$ as above.
Assume that the $Q_{9r_j}(c_j)$ are pairwise disjoint and miss $\Gamma$.
Give integer weights $w_j$ and a rational triangle chain $T$ satisfying
the finite oriented-boundary identity
\[
 \partial T=\Gamma-\sum_j w_j\partial Q_{r_j}(c_j).
\]
Subdivisions at common rational vertices are allowed. Require
$\lVert z-c_j\rVert_\infty\ge r_j/16$ throughout the carrier of $T$,
and give the holomorphic chart data for $f$ on that carrier and
the shells $S(c_j,r_j)$, with compatible values on overlaps.
Distances and domains here refer to rational cells. Such a filling for
a rectangle with square holes can be produced by drawing the horizontal
and vertical side lines and splitting each retained rectangle into
two triangles. Holomorphy on the carrier of the filling is essential.

\paragraph{Theorem.}
For each $j$, the translated local chart constructs the Laurent
coefficient $R_j$ of degree $-1$ at $a_j$. With name equality
interpreted at every requested precision,
\[
 \int_\Gamma f(z)\,dz=2\pi i\sum_j w_j R_j,\qquad
 w_j=(2\pi i)^{-1}\int_\Gamma \frac{dz}{z-a_j}.
\]
Here $R_j$ is the coefficient of the presented annulus centred at $a_j$.
When the data extend to a punctured neighbourhood of $a_j$ as specified
below, it is the residue of that point. This includes isolated essential
singularities.

\paragraph{Proof.}
Fix $j$ and suppress its index. For $n\ge b$, join the four vertices
of $Q_r(c)$ to the corresponding translated vertices of $Q_r(A(n))$.
Each of the four quadrilaterals is a signed pair of rational triangles.
Its points are of the form $c+x+t(A(n)-c)$ with $x$ on the corresponding
edge of $Q_r(0)$ and $0\le t\le1$. Their distances from $c$ lie between
$r-r/32$ and $r+r/32$ in the sup norm. Thus Cauchy's triangle theorem
applies there, including to degenerate panels. The interior diagonals
and the four connector edges cancel with opposite signs, giving the
exact equality of the integrals on the two rational squares.

The difference between the integral over $\partial Q_r(A(n))$ and
$\int_{\partial Q_r(0)}g(x)\,dx$ has norm at most $8r2^{-u}$ whenever
$n\ge\max(b,1+\nu(u))$. Indeed its two integrands in the fixed $x$
coordinate are within $2^{-u}$ at every rational quadrature node by
the common-later-centre argument. The increments in each finite edge
sum coincide and their total norm is $8r$. The quadrature error on
either side has its supplied modulus and can be made arbitrarily small.
This proves equality for the translated square. By the translation lemma
and the rational annular theorem, its integral is $2\pi iR_j$.
Applying Cauchy to $T$ and using the boundary identity now gives the
first formula.

For the second formula apply the same chain calculation to
$k_j(z)=(z-a_j)^{-1}$. On $T$ and its own shell the distance to $a_j$
is at least $r_j/16-r_j/32$. The disjoint outer squares keep
$k_j$ holomorphic throughout every other local square. The inverse
evaluator is effective here. When the sup distance has lower bound
$\delta>0$, take centre error at most $\delta/2$, use
$|v^{-1}|_1\le2/\lVert v\rVert_\infty$ and the inverse difference
identity to budget its value error. Squaring gives derivative
$-v^{-2}$. The identity
$(v+h)^{-1}-v^{-1}+h/v^2=h^2/(v^2(v+h))$
gives a remainder constant $U^3$ for any common endpoint inverse bound
$U$. In the translated chart at $a_j$ this kernel is $1/x$ and has
residue one. Its integral around each other square vanishes by its
two-triangle filling. The second formula follows. In particular changing
a rational filling or its subdivision gives the same integers since
these normalized integral names depend only on $a_j$ and $\Gamma$.
Independence of radius follows by Cauchy on the intervening annulus
in translated coordinates.\hfill$\square$

\paragraph{From an annular coefficient to a point residue.}
A single shell can surround several singularities. For example,
$(z-d)^{-1}$ with $0<|d-a|_1<r/64$ contributes one to the
annular coefficient around $a$, although the function is regular at $a$.
A sufficient punctured-germ presentation consists of the preceding shell
data for every $r_m=2^{-m}r$, using $b_m=b+m$ and $c_m=A(b_m)$,
jointly PR in $m$ and agreeing on overlaps. In translated coordinates
the adjacent evaluations agree on the common annulus by the
common-later-centre calculation. The square annulus between
half-sides $r_m/2$ and $r_m$ lies in that common annulus, so Cauchy's
theorem on its rectangular filling identifies successive coefficient
integrals. The common value is the coefficient at the isolated point.
For a numerical depth $m$ and precision $q$, sum the $m$ adjacent
differences at precision $q+1+\sigma(m+1)$. This proves independence of
depth with a rational prefix invariant. The positive gap around each
other centre is retained at all the smaller scales.

\paragraph{Rectifiable curves.}
The same result holds for a closed PR-named curve $\gamma$ with
a PR continuity modulus $\nu_\gamma$ and a supplied upper bound on the
sum of its increments' norms over every dyadic partition. Give a
positive rational $\rho$ such that $f$ has the chart data on the
$\rho$-collar of the curve and
$\lVert\gamma(t)-a_j\rVert_\infty\ge\rho$ for every $j,t$.
Choose $2^{-C}\le\rho/8$, $m=\nu_\gamma(C)$, and join
$2^{-C}$-accurate rational samples at the times $i/2^m$ to obtain a
rational polygon $P$, using the same sample at both endpoints. On each
subinterval both endpoint samples are within $2\cdot2^{-C}$ of
$\gamma(t)$, so $|P(t)-\gamma(t)|_1\le\rho/4$ by convexity.
The linear homotopy from $\gamma$ to $P$ has its $\rho/2$-collar
in the chart and at sup distance at least $\rho/4$ from all centres.
For a PR modulus use $\max(\nu_\gamma(k+2),k+2+d)$ where $2^d$ bounds
one, $\rho/4$, and the greatest speed on the polygon's parameter
segments. Changing $t$ and then the homotopy parameter in its affine
formula proves this bound. The rational-grid homotopy theorem in the
linked report identifies the two contour integrals, for $f$ and for
each inverse kernel individually. Whenever $P$ has the filling and
shell certificates above, the same coefficients and integers therefore
give the value and the windings for $\gamma$.

\paragraph{Precision and induction.}
For a target $2^{-q}$ take
$u=q+4+\sigma(1+\sum_j8r_j|w_j|)$. The total weighted error of replacing
moving contours by frozen samples is then at most $2^{-q-4}$.
A local Laurent truncation of order
$N\ge1+q+4+\sigma(1+\sum_j 8r_j|w_j|40M_j)$ gives the same
bound on the accumulated truncation error. For the residual of Cauchy's theorem on $T$ and on all rational
translation panels, give each constituent triangle error at most
$2^{-q-4-\sigma(1+L)}$ in the upstream triangle schedule, where $L$
is the sum of absolute triangle multiplicities in $T$ plus
$8\sum_j|w_j|$ for the panels. This bounds their combined residual by
$2^{-q-4}$. Finally set $W=\sum_j|w_j|$ and
$B_*=\sum_j |w_j|8r_jM_j$. The local coefficient integral and
$|(2\pi i)^{-1}|_1\le1$ give $\sum_j|w_j|\,|R_j|_1\le B_*$.
If the sampled $\pi$ has magnitude at most five, and its error and the
errors in the outer integral and individual residues are
$\eta_\pi,\eta_\Gamma,\eta_R$, the additional error in the equality is
at most $\eta_\Gamma+2B_*\eta_\pi+10W\eta_R$. Use $q+4$ bits for the
outer integral, $q+4+\sigma(1+2B_*)$ for $\pi$, and
$q+4+\sigma(1+10W)$ for the residue names, with the rational-case
schedules for quadrature and numerical coefficient evaluation. Even
adding all six indicated error allowances gives less than $2^{-q}$.
The proof of cancellation fixes an elementary oriented edge and inducts
on a partial sum of signed incidences. The sum estimate inducts on the
rational accumulator of the norm bounds at fixed precision. The
analytic hypotheses are used by instantiation at the rational sample
points. Equality of names means agreement to all rational accuracies.
A violated rational inequality has a decidable positive gap and a PR
precision that detects it.

\paragraph{Example.}
Take $a=\sqrt2$, $f(z)=(z^2-2)^{-1}$, and let $\Gamma$ be the positive
boundary of $[1,2]\times[-1/2,1/2]$. Use the lower dyadic name
$A(n)=2^{-n}\lfloor\sqrt{2\cdot 4^n}\rfloor$, obtained by bounded
integer search, and put $r=1/64$, $b=12$, $c=181/128$.
The outer square of half-side $9r$ lies strictly inside the rectangle.
The four lines supporting the hole divide the rectangle into nine cells.
Deleting the middle cell gives a filling of sixteen triangles with
weight one. On the carrier and shell the factor $z-a$ is bounded away
from zero by the clearance above and $\lVert z+a\rVert_\infty\ge1$
by its real part. The partial-fraction identity
\[
 \frac1{z^2-2}=\frac1{2a}\left(\frac1{z-a}-\frac1{z+a}\right)
\]
supplies effective data by the inverse estimate, and gives
$R=1/(2a)$ because the second summand is holomorphic across the
whole square. Thus $\int_\Gamma dz/(z^2-2)=\pi i/\sqrt2$.
The arithmetic fact $a^2=2$ follows directly from
$A(n)^2\le2<(A(n)+2^{-n})^2$ and $A(n)\le2$.

\paragraph{Lean implementation path.}
The host file \texttt{ComplexCase.lean} defines an incidence functional
on ordered vertex pairs, the boundary of one oriented triangle, and
the signed two-triangle panel. Its lemmas \texttt{strip\_boundary}
and \texttt{closed\_strip} give cancellation for every finite strip,
and \texttt{annulus\_clearance} proves the integer-scaled shell
inequality. The file also proves both taxicab bounds for Gaussian integer
multiplication. Its \texttt{derivative\_variation} lemma recovers the
factor four in the derivative modulus after clearing positive denominators. If $|Ch|_1\le 2E|h|_1$ and $|h|_1>0$, then $|C|_1\le4E$. The companion \texttt{ComplexObject.lean} proves the strip identity in
the repository's object calculus. For an arbitrary externally fixed
path and an incidence table $E$ whose entries are signed pairs of open
terms, \texttt{ResidueObject.closed\_strip} constructs a
$\Gamma\vdash_{\mathsf{PRAAx}}$ derivation of the signed boundary
equation. Its \texttt{square\_strip} specialization gives the eight
triangles used above. The proof interprets the incidence expressions
in the repository's ring of formal differences, cancels them there, and
applies \texttt{ZT.derive}. An implementation of the geometry must still
supply the incidence terms for the chosen rational test edge, with
normalization of endpoints and signed comparisons. A single uniform
object formula for coded paths of variable length would additionally
use a quantifier-free prefix induction. Next implement
$|\cdot|_1,\lVert\cdot\rVert_\infty$,
the rational inverse and the two displayed scheduling functions by
cross-multiplication and bounded recursion. The derivative lemma uses
the two instantiated remainder inequalities and the inequality
$(|x|+|y|)^2\le2(x^2+y^2)$.
The translated samples are the displayed compositions of PR functions.
With these arithmetic lemmas the remaining subgoals are the
rational triangle and Laurent derivations from the upstream roadmap,
followed by universal elimination at the translated programs and
the finite weighted sum. The included Python example checks the grid
boundary, all triangle-strip incidences, dyadic brackets, and the
rational-proxy partial-fraction identity exactly.
